\section{Puente 1: Transformación de los datos de entrenamiento con In-Context Augmentation}

La primera ruta de puente se completa antes del entrenamiento: dado que el modelo puede observar reversal, missing link y multi-hop implication en el context, primero haga explícitas estas relaciones dentro del context y luego fine-tune con los datos ampliados. Esta estrategia amortiza un proceso de inferencia costoso fuera de línea hacia los parámetros, siendo adecuada para escenarios donde las familias de problemas futuros sean relativamente predecibles.

\subsection{Offline vs Online: división arquitectónica primero}

La diapositiva 25 divide la propuesta en dos categorías. Los métodos Train-time/offline generan conclusiones complementarias antes del entrenamiento, haciendo que la inferencia posterior sea más barata; los métodos test-time/online conservan las experiencias originales y recuperan o regeneran información cuando surgen las preguntas, cubriendo un espectro más abierto pero siendo más costosos por solicitud. Esta división es más importante que los nombres específicos de los algoritmos, ya que determina la gobernanza de datos, el latencia, el almacenamiento en caché, la evaluación y los métodos de recuperación ante fallos.

\lecturefigure{slide-025.jpg}{Las dos grandes familias para cerrar la brecha latent: augmentation offline y recuperación online}{Diapositiva oficial p. 25 del curso Stanford CS25 V6 Lecture 06}

\begin{importantbox}{Pregunte primero si los futuros problemas son enumerables}
Si los tipos de consulta futuras son estables, el entrenamiento se puede reejecutar y la latencia de inferencia es sensible, la augmentation offline tiene atractivo; si los objetivos cambian constantemente, los detalles de cada experiencia individual son importantes o se requiere rastrear evidencia original ante errores, entonces la recuperación/regeneración online es más adecuada.
\end{importantbox}

\subsection{Algoritmo: conécte primero los documentos y luego fine-tune}

La in-context augmentation coloca el corpus completo o fragmentado en el context, exigiendo al modelo reescribir el documento actual, conectar otros documentos y completar relaciones inversas e implicaciones transversales a documentos. Los resultados generados se combinan con los datos originales para proceder al fine-tuning. La clave aquí no es un simple paraphrase, sino permitir que las habilidades ICL existentes conviertan la latent relation en una señal de supervisión directa.

\lecturefigure{slide-026.jpg}{Enfoque 1: el pipeline de entrenamiento de in-context augmentation}{Diapositiva oficial p. 26 del curso Stanford CS25 V6 Lecture 06}

\begin{lstlisting}[language=Python,caption={Marco de implementación mínima para in-context augmentation}]
augmented = []
for document in corpus:
    context = retrieve_related_documents(corpus, document)
    prompt = build_connection_prompt(context, document)
    implications = teacher_model.generate(prompt)
    augmented.extend(validate_and_deduplicate(implications))

finetune_dataset = corpus + augmented
student = finetune(base_model, finetune_dataset)
\end{lstlisting}

\lecturefigure{slide-027.jpg}{Conectar un documento individual con el corpus global para generar conclusiones entrenables directamente}{Diapositiva oficial p. 27 del curso Stanford CS25 V6 Lecture 06}

Al analizar este ejemplo, rastree tres capas de objetos: los documentos originales solo escriben hechos locales; el context global proporciona otra relación; y el modelo de augmentation genera conclusiones transversales a documentos como "Dax can zarp". El estudiante final no necesita escanear todo el corpus nuevamente durante la prueba, ya que dicha conclusión se ha convertido en un objetivo de entrenamiento.

\subsection{Resultados: augmented fine-tuning puede igualar o superar ICL}

En familias de problemas como reversal y syllogism que ya han sido cubiertas por augmentation, el fine-tuning con datos ampliados puede igualar al ICL, siendo incluso mejor en algunos configuraciones. Esto no es misterioso: los resultados de razonamiento del ICL se convierten en más supervisión directa, otorgando al estudiante la dirección de objetivos que faltaba en el fine-tuning original.

\lecturefigure{slide-028.jpg}{Augmented fine-tuning iguala o supera al ICL en reversal y syllogism}{Diapositiva oficial p. 28 del curso Stanford CS25 V6 Lecture 06}

\lecturefigure{slide-029.jpg}{Conclusión: el ICL puede funcionar como un expansor de relaciones para los datos de entrenamiento}{Diapositiva oficial p. 29 del curso Stanford CS25 V6 Lecture 06}

En la grabación de audio [00:24:17--00:27:21], el docente denomina a ambos resultados como "muy buenos pero diferentes": el ICL realiza un razonamiento flexible directamente durante la prueba; la augmentation genera conexiones con el ICL primero y luego destila esas conexiones en los parámetros. La ventaja de este último solo se aplica a las familias de problemas que han sido expandidas.

\begin{knowledgebox}{Por qué augmented fine-tuning puede superar al ICL directo}
El ICL directo debe recuperar y combinar información en un context largo cada vez, sufriendo de efectos de posición, interferencia y capacidades de recuperación del context; la augmentation puede generar múltiples direcciones, deduplicar y filtrar, y luego reforzar la misma conclusión mediante optimización multi-etapa. Por tanto, "superar al ICL" no indica que el aprendizaje de parámetros sea inherentemente más flexible, sino que la pipeline offline ha construido una distribución de supervisión más directa para la tarea objetivo.
\end{knowledgebox}

\subsection{No desde cero: lo que cambia es la accesibilidad}

La diapositiva 30 responde anticipadamente a las dudas de teoría de la información: si la augmentation solo lee los datos originales, ¿cómo puede "obtener más información"? La respuesta correcta es que no crea nueva ground-truth; utiliza el conocimiento previo existente del modelo y el razonamiento del context para reescribir las conclusiones implícitas en los datos originales como muestras explícitas. Al comprender estas dos páginas, debe distinguir entre semantic implication y independent evidence: la primera puede ser expandida mediante cálculos, mientras que la segunda sigue debiendo provenir de observaciones reales o fuentes confiables; la augmentation solo puede mejorar la representación y el acceso, no reemplazar la recolección de evidencia.

\lecturefigure{slide-030.jpg}{Pregunta: ¿la augmentation genera información desde nothing?}{Diapositiva oficial p. 30 del curso Stanford CS25 V6 Lecture 06}

\lecturefigure{slide-031.jpg}{Respuesta: convertir la latent implication en supervisión explícita}{Diapositiva oficial p. 31 del curso Stanford CS25 V6 Lecture 06}

Si el amplificador $A$ es determinístico para los datos $D$, entonces
$$
H\!\left(A(D)\mid D\right)=0.
$$
Donde:
\begin{itemize}
  \item $A(D)$ son las reescrituras, relaciones inversas o conclusiones combinadas derivadas de $D$;
  \item $H(\cdot\mid\cdot)$ es la entropía condicional;
  \item Esta fórmula indica que, dado los datos originales, una augmentation determinística no agrega información aleatoria independiente nueva;
  \item La generación real de LLMs es estocástica e inyecta el conocimiento previo del modelo y errores, por lo tanto se debe verificar la factuality y no tratar todos los nuevos tokens como nueva evidencia.
\end{itemize}

\begin{warningbox}{“La información ya está en los datos” no equivale a “la generación es necesariamente correcta”}
Las conclusiones válidas de Syllogism y las conexiones hallucinadas pueden ser escritas fluidamente por el modelo. Las pipelines de producción deben registrar los documentos fuente, tipos de razonamiento, modelos generadores y versiones del prompt, y filtrar usando reglas, verificadoros, muestreo cruzado o inspección manual; de lo contrario, la expansión latent solidificará errores también dentro de los parámetros.
\end{warningbox}

\subsection{Límite fundamental: no se puede enumerar con anticipación cada futuro uso}

Los métodos offline deben adivinar qué preguntará el futuro. Un investigador que lee un artículo hoy no puede escribir con anticipación el significado de esa experiencia para todos los temas, objetivos y combinaciones futuros; el número de relaciones crece con las combinaciones de documentos y objetivos potenciales. La diapositiva 32 ilustra esto con el ejemplo de "quizás una tesis doctoral futura lo use", explicando que la augmentation completa es inviable en un mundo abierto. Un límite más profundo es que los propios objetivos cambian: la forma disponible de un mismo hecho para recuperación, planificación, explicación, contrafactual y auditoría de seguridad no es idéntica, por lo tanto generar con anticipación varios paraphrases fijos no puede cubrir las necesidades computacionales futuras.

\lecturefigure{slide-032.jpg}{Pre-aumentar datos para cada objetivo futuro no es escalable}{Diapositiva oficial p. 32 del curso Stanford CS25 V6 Lecture 06}

Asumiendo que el corpus tiene $N$ experiencias y cada inferencia combina a lo sumo $k$ experiencias, el número potencial de combinaciones alcanza aproximadamente
$$
\sum_{i=1}^{k}\binom{N}{i},
$$
donde $N$ es el número de experiencias y $k$ la profundidad de combinación permitida. Los problemas reales también incluyen objetivos, idioma y formato de salida, ampliando aún más el espacio. Por tanto, la augmentation offline requiere conocimiento previo de tareas y presupuesto, no una expansión infinita.

\subsection{Resumen de la sección}

El valor de la in-context augmentation reside en transformar las latent relation que se forman fácilmente dentro del context en supervisión explícita, y trasladar el cálculo durante la prueba al momento del entrenamiento. Es adecuada para familias de problemas predecibles, pero asume costos de generación, verificación, inflación de datos y deriva del modelo; el límite más importante es la incapacidad de prever todos los usos futuros de cada experiencia. Esto es precisamente el punto de partida para la recuperación episódica online.
